%! Author = dario
%! Date = 2/7/25

% Preamble
\documentclass[11pt]{article}

% Packages
\usepackage{amsmath}



\title{Msc thesis}
\date{2025}

\usepackage{biblatex}
\usepackage{hyperref}
\addbibresource{bib.bib}

%\author{Dario Filatrella}
% write authors and supervisors
\title{M.Sc. Thesis Draft}
\author{Dario Filatrella \\ Supervisors: Aida Khajavirand and Antonio De Rosa}
\date{2025}

% Document
\begin{document}

    \maketitle



    \section*{Preface}
    This document presents preliminary research on branched transportation, conducted under the supervision of Professors
    De Rosa and Khajavirad. It is currently structed as a draft publication and I will later adapt it to the thesis format.
    On GitHub, at the link \url{https://github.com/DarioFi/optimal_transport} you can find the repository containing all the work done.
    It includes source code for result generation, raw experimental data, and Jupyter notebooks containing
    experimental results and visualizations not featured in the main text.


    \section*{Abstract}
    This work investigates the discrete
    branched transportation problem as described in~\cite[Bernot, Caselles, Morales]{1}.
    Our research is divided into two main parts.
    First, we establish a sharper bound on the angles between edges in optimal solutions, leading to an improved
    bound on the degree of the nodes in the optimal configuration for a given mass distribution and cost parameter $\alpha$.

    Second, we analyze a mixed-integer nonlinear programming (MINLP) formulation of the problem.
    While demonstrating that this MINLP approach successfully yields globally optimal solutions, we show that there are significant
    performance limitations and opportunities for improvement.

    We provide a detailed analysis of the key factors affecting the performance of the BARON solver in this context,
    offering insights for future algorithmic developments.

    Our work on finding globally optimal solutions addresses an underexplored aspect of branched transportation.
    This approach yields valuable insights into both the structural properties of optimal solutions and
    the fundamental behavior of the problem, contributing to both theoretical understanding and computational methodology.


\end{document}